% !TeX root = ../ADCS & GNC Documentation.tex
\chapter{State Estimation}

\section{Overview}

Attitude estimation is performed using a Multiplicative Extended Kalman Filter (MEKF).

\section{Asynchronous Sensor Fusion}

The state estimation architecture is designed to operate with asynchronous sensor updates. This is required due to the differing update rates of the onboard sensors.

The IMU provides angular velocity measurements at a high rate, while the star tracker provides attitude measurements at a significantly lower rate (on the order of 2 Hz). Additionally, due to uncertainty in the CAN bus communication protocols, exact timings of these sensor updates cannot be guaranteed.

The Multiplicative Extended Kalman Filter (MEKF) is structured to accommodate this disparity in measurement availability:

\begin{itemize}
	\item IMU measurements are used continuously in the propagation step
	\item Star tracker measurements are incorporated only when available for correction
\end{itemize}

Between measurement updates, the filter performs state propagation using a zero-order hold assumption on angular velocity. When higher-rate IMU data is available, midpoint integration is used to improve propagation accuracy.

This asynchronous design ensures that the estimator remains stable and accurate despite the differing sensor update rates.

\section{State Definition}

\[
x =
\begin{bmatrix}
	\delta \theta \\
	b
\end{bmatrix}
\]

where $\delta \theta$ represents the small-angle attitude error and $b$ represents gyro bias.

\section{Prediction}

\begin{itemize}
	\item Quaternion propagation via exponential map
	\item Covariance propagation using linearized dynamics
\end{itemize}

\section{Correction}

\begin{itemize}
	\item Innovation derived from quaternion error
	\item Joseph-form covariance update for numerical stability
\end{itemize}

\section{Implementation Details}

\begin{itemize}
	\item Hemisphere enforcement is required for stability
	\item Midpoint integration is used when IMU updates are available
	\item Zero-order hold is used otherwise
\end{itemize}